\section{Post-Training: RLHF y Constitutional AI}

El preentrenamiento aporta un conocimiento amplio del lenguaje y del mundo, pero no produce por sí solo una conducta conversacional acorde con la intención del producto. El Post-training incorpora al modelo instruction, preference, principle y tool policy. La clase avanza de RLHF a Constitutional AI/RLAIF; lo importante no es sustituir por completo el human feedback, sino ampliar la escala de la supervisión y lograr que los objetivos de conducta sean más explícitos e iterables.

\subsection{RLHF: convertir preferencias en una proxy reward mediante entrenamiento}

El \term{reinforcement learning from human feedback} (RLHF) suele incluir supervised fine-tuning, human preference comparison, un \term{reward model} (modelo de recompensa) y policy optimization. El reward model aprende a predecir qué respuesta prefieren las personas; después, el RL optimiza el modelo para aumentar la predicted reward, con restricciones como la KL para evitar que la política se aleje demasiado.

\lecturefigure{09-rlhf-pipeline.png}{RLHF encadena demonstrations, preference, reward model y policy optimization en un pipeline de entrenamiento.}{InstructGPT; subtítulos locales 21:00--26:20; reconstrucción conceptual.}

\begin{knowledgebox}{Lectura de la figura: cada capa tiene su propio error}
El SFT está limitado por la cobertura de las demonstrations; la preference depende del annotator y de la rubric; el reward model puede ser objeto de exploit ante un distribution shift; la RL policy puede elevar la proxy reward y, al mismo tiempo, perjudicar la helpfulness real. Al final, la human evaluation, la safety evaluation y las product metrics no pueden sustituirse por la reward curve.
\end{knowledgebox}

\begin{warningbox}{El reward hacking es un problema de proxy, no solo un bug del algoritmo}
Ninguna rubric finita puede expresar por completo «útil, honesto, inofensivo y adecuado al contexto». El modelo puede aprender a complacer al judge, a negarse en exceso o a recurrir a un estilo superficial para obtener una puntuación alta. Conviene usar evaluation multidimensional, adversarial examples, holdout tasks y policy review, para evitar que una sola reward se convierta en la única verdad.
\end{warningbox}

\subsection{Constitutional AI y RLAIF}

\term{Constitutional AI} usa un conjunto de principios para guiar al modelo en la critique y la revise de sus propias respuestas, y con ello genera AI preference data; el \term{reinforcement learning from AI feedback} (RLAIF) usa un evaluator model para producir preferencias o una reward signal. Las personas siguen encargándose de elegir los principles, auditar la conducta, resolver conflictos y decidir la deployment boundary; el AI feedback se encarga de reducir el costo de que cada muestra requiera anotación humana.

\lecturefigure{10-constitutional-ai.png}{Constitutional AI amplía la supervisión mediante principles, critique/revision, AI preferences y RLAIF.}{Anthropic Constitutional AI; subtítulos locales 26:00--27:40; reconstrucción conceptual.}

\begin{knowledgebox}{Lectura de la figura: la capacidad del base model es un requisito previo}
CAI solo funciona si el modelo es capaz de entender los principios, detectar problemas y proponer una revision mejor; además, el evaluator model puede compartir los blind spots del target model. Los conflictos entre principios, el contexto cultural, la negativa excesiva y el hidden reasoning requieren human audit. RLAIF amplía el volumen de supervisión, pero eso no equivale a eliminar la human governance.
\end{knowledgebox}

\teachervoice{Mann describe Constitutional AI como un avance que llegó tras cruzar un umbral de capacidad: los primeros modelos no eran capaces de hacer una critique confiable de sí mismos, y solo cuando los modelos se volvieron más fuertes pudieron convertirse en herramientas de supervisión. Este ejemplo muestra que el capability growth también puede crear nuevas safety techniques, y no solo aumentar los riesgos.}

\subsection{Resumen de la sección}

El Post-training es una optimización de proxies en varias capas. Tanto RLHF como RLAIF requieren principios claros, evaluation independiente y human governance; modifican la conducta, pero no sustituyen la base capability, el system control ni la deployment safety.
